% ============ 19 有限性定理; Z_ℓ 模层 (原书 pp.123-126) ============
\chapter{\texorpdfstring{有限性定理; $\mathbb{Z}_{\ell}$ 模层}{有限性定理；Z\_l 模层}}\label{sec:19}
\index{有限性定理}
\index{$\mathbb{Z}_{\ell}$ 模层}
\index{$\ell$ 进层}
\index{$\mathbb{Q}_{\ell}$ 模层}
\index{lisse 层}

\subsection*{有限性定理}

\begin{theorem}\label{thm:19.1}
设 $X$ 是可分闭域 $k$ 上的簇，$F$ 是 $X_{\mathrm{et}}$ 上的可构造层.在如下两种情形中，群 $H^{r}(X_{\mathrm{et}}, F)$ 都有限.
\begin{enumerate}
  \item $X$ 完备，或
  \item $F$ 没有 $p$-挠，其中 $p$ 是 $k$ 的特征.
\end{enumerate}
\end{theorem}

\begin{proof}
情形 (a) 是紧合基变换定理的一部分，在第 17 节讨论.

对非奇异簇 $X$ 与 $X_{\mathrm{et}}$ 上的局部常值层 $F$，可以通过对 $X$ 的维数作归纳并利用初等纤维化(\S 21)的存在性证明 $H^{r}(X_{\mathrm{et}}, F)$ 有限---见 SGA 4, VI.5.2.

或者，在这种情况下，它由 Poincar\'e 对偶定理(见后面)推出，后者表明
\[
H^{r}(X, F) \text{ 与 } H^{2m-r}_{c}(X, \mathcal{L}F(m)) \text{ 对偶}, \qquad m = \dim X.
\]

一般情形更困难---见 SGA $4\frac{1}{2}$, pp.233--261.[但当然，用 de Jong 的分解定理可以简化证明.]
\end{proof}

\subsection*{\texorpdfstring{$\mathbb{Z}_{\ell}$ 模层}{Z\_l 模层}}

到目前为止，我们只谈论了挠层.然而，对我们而言，拥有特征零域上的向量空间作为上同调群很重要，例如为了有好的 Lefschetz 不动点公式.然而，非挠层系数的平展上同调群是异常的.例如，当 $X$ 正规时，
\[
H^{1}(X_{\mathrm{et}}, \mathbb{Z}) = \operatorname{Hom}_{\mathrm{conts}}(\pi_1(X, \bar{x}), \mathbb{Z}),
\]
($\mathbb{Z}$ 带离散拓扑)，它为零，因为连续同态 $f \colon \pi_1(X, \bar{x}) \to \mathbb{Z}$ 在 $\pi_1(X, \bar{x})$ 的某个开子群上必须为零，而这样的子群有有限指数($\pi_1(X, \bar{x})$ 紧).\footnote{更好的是，假定 $X$ 正规，设 $i \colon x \to X$ 是一般点的包含.则 $\mathbb{Z}_X \cong i_{*}\mathbb{Z}$，且对 $j = 1$，$R^{j}i_{*}\mathbb{Z} = 0$，从而 $H^{1}(X_{\mathrm{et}}, \mathbb{Z}_X) \cong H^{1}(X_{\mathrm{et}}, i_{*}\mathbb{Z}) \cong H^{1}(x_{\mathrm{et}}, \mathbb{Z})$，它为零.当 $X$ 是以单个结点为奇异点的连通曲线时，$H^{1}(X_{\mathrm{et}}, \mathbb{Z}) \cong \mathbb{Z}$.要看出这一点，考察正合序列 $0 \to \mathbb{Z} \to \varphi_{*}\mathbb{Z} \to \varphi_{*}\mathbb{Z}/\mathbb{Z} \to 0$ 的上同调，其中 $\varphi \colon \widetilde{X} \to X$ 是 $X$ 的正规化.}类似地，按显然的定义，$H^{1}(X_{\mathrm{et}}, \mathbb{Z}_{\ell}) = 0$(它由对 $\mathbb{Z}_{\ell}$ 的离散拓扑连续的同态 $\pi_1(X, \bar{x}) \to \mathbb{Z}_{\ell}$ 组成).解决办法是定义
\[
H^{r}(X_{\mathrm{et}}, \mathbb{Z}_{\ell}) = \varprojlim_{n} H^{r}(X_{\mathrm{et}}, \mathbb{Z}/\ell^{n}\mathbb{Z})
\]
---上同调不与层的逆向极限交换.按这个定义，
\[
H^{1}(X_{\mathrm{et}}, \mathbb{Z}_{\ell}) \overset{\mathrm{def}}{=} \varprojlim_{n} H^{1}(X_{\mathrm{et}}, \mathbb{Z}/\ell^{n}\mathbb{Z}) \cong \varprojlim_{n} \operatorname{Hom}_{\mathrm{conts}}(\pi_1(X, \bar{x}), \mathbb{Z}/\ell^{n}\mathbb{Z})
\]
\[
\cong \operatorname{Hom}_{\mathrm{conts}}(\pi_1(X, \bar{x}), \mathbb{Z}_{\ell}),
\]
其中 $\mathbb{Z}/\ell^{n}\mathbb{Z}$ 带离散拓扑，$\mathbb{Z}_{\ell}$ 带 $\ell$ 进拓扑(读者应检验最后的 $\cong$).

给出有限生成 $\mathbb{Z}_{\ell}$-模 $M$ 与给出族 $(M_n, f_{n+1} \colon M_{n+1} \to M_n)_{n \in \mathbb{N}}$ 是一回事，后者满足：
\begin{enumerate}
  \item 对所有 $n$，$M_n$ 是有限 $\mathbb{Z}/\ell^{n}\mathbb{Z}$-模；
  \item 对所有 $n$，映射 $f_{n+1} \colon M_{n+1} \to M_n$ 诱导同构 $M_{n+1}/\ell^{n}M_{n+1} \to M_n$.
\end{enumerate}
给定 $M$，取 $M_n = M/\ell^{n}M$，$f_{n+1}$ 为商映射.反过来，给定 $(M_n, f_{n+1})$，定义 $M = \varprojlim_{n} M_n$.事实上，对应 $M \leftrightarrow (M_n, f_{n+1})$ 可以做成范畴的等价\footnote{设
\[
M_n = \mathbb{Z}/\ell^{n}\mathbb{Z} \times \mathbb{Z}/\ell^{n-1}\mathbb{Z} \times \cdots \times \mathbb{Z}/\ell\mathbb{Z},
\]
$M_{n+1} \to M_n$ 是显然的商映射.则 $M = \varprojlim M_n = \prod \mathbb{Z}_{\ell}$(由正整数指标化的 $\mathbb{Z}_{\ell}$ 的副本的乘积).这个例子表明有限 $\mathbb{Z}_{\ell}$-模的逆向极限不必是有限生成 $\mathbb{Z}_{\ell}$-模.然而，这个例子不满足 $M_n = M_{n+1}/\ell^{n}M_{n+1}$ 的条件，因此与本段中的论断不矛盾.不过，也许应该加条件：对某个 $n_0$，系统 $N_n \overset{\mathrm{def}}{=} \ell^{n_0}M_{n+n_0}$ 平坦.由有限生成 $\mathbb{Z}_{\ell}$-模 $M$ 得到的任何系统 $M/\ell^{n}M$ 都满足这个条件，其中 $n_0$ 选取为使 $\ell^{n_0}$ 消灭 $M$ 的挠的数.}.

设 $(M_n, f_n)_{n \in \mathbb{N}}$ 是满足 (a),(b) 的系统.由归纳，得到典范同构 $M_{n+s}/\ell^{n}M_{n+s} \cong M_n$.把
\[
0 \to \mathbb{Z}/\ell^{s}\mathbb{Z} \xrightarrow{\text{``}\ell^{n}\text{''}} \mathbb{Z}/\ell^{n+s}\mathbb{Z} \to \mathbb{Z}/\ell^{n}\mathbb{Z} \to 0
\]
与 $M$ 作张量积，得到序列
\[
0 \to M_s \to M_{n+s} \to M_n \to 0,
\]
若 $M$ 平坦(等价地，无挠)，则它正合.

上面的讨论启发如下定义：$X$ 上的 $\mathbb{Z}_{\ell}$-模层(或 $\ell$ 进层)是族 $(M_n, f_{n+1} \colon M_{n+1} \to M_n)$，使得
\begin{enumerate}
  \item 对每个 $n$，$M_n$ 是 $\mathbb{Z}/\ell^{n}\mathbb{Z}$-模的可构造层；
  \item 对每个 $n$，映射 $f_{n+1} \colon M_{n+1} \to M_n$ 诱导同构 $M_{n+1}/\ell^{n}M_{n+1} \to M_n$.
\end{enumerate}

设 $(M_n, f_n)_{n \in \mathbb{N}}$ 是 $X$ 上的 $\mathbb{Z}_{\ell}$-模层.由归纳，得到典范同构 $M_{n+s}/\ell^{n}M_{n+s} \cong M_n$.把
\[
0 \to \mathbb{Z}/\ell^{s}\mathbb{Z} \xrightarrow{\text{``}\ell^{n}\text{''}} \mathbb{Z}/\ell^{n+s}\mathbb{Z} \to \mathbb{Z}/\ell^{n}\mathbb{Z} \to 0
\]
与 $M_{n+s}$ 作张量积，得到序列
\[
0 \to M_s \to M_{n+s} \to M_n \to 0.
\]
如果对所有 $n$ 与 $s$ 这个序列都正合，我们说 $M$ 平坦.

对 $\mathbb{Z}_{\ell}$-模层 $M = (M_n)$，定义
\[
H^{r}(X_{\mathrm{et}}, M) = \varprojlim_{n} H^{r}(X_{\mathrm{et}}, M_n), \qquad H^{r}_{c}(X_{\mathrm{et}}, M) = \varprojlim_{n} H^{r}_{c}(X_{\mathrm{et}}, M_n).
\]
例如，若用 $M_n$ 为常值层 $\mathbb{Z}/\ell^{n}\mathbb{Z}$，$f_n$ 显然的 $\mathbb{Z}_{\ell}$-模层记 $\mathbb{Z}_{\ell}$，则
\[
H^{r}(X_{\mathrm{et}}, \mathbb{Z}_{\ell}) \overset{\mathrm{def}}{=} \varprojlim H^{r}(X_{\mathrm{et}}, \mathbb{Z}/\ell^{n}\mathbb{Z}).
\]

\begin{theorem}\label{thm:19.2}
设 $M = (M_n)$ 是域 $k$ 上的簇 $X$ 上的平坦 $\mathbb{Z}_{\ell}$-模层.假定 $k$ 可分闭，且 $X$ 完备或 $\ell \neq \operatorname{char}(k)$.则每个 $H^{r}(X_{\mathrm{et}}, M)$ 都有限生成，且有上同调群的正合序列
\[
\cdots \to H^{r}(X_{\mathrm{et}}, M) \xrightarrow{\ell^{n}} H^{r}(X_{\mathrm{et}}, M) \to H^{r}(X_{\mathrm{et}}, M_n) \to H^{r+1}(X_{\mathrm{et}}, M) \to \cdots.
\]
\end{theorem}

\begin{proof}
对每个 $s \geq 0$，得到正合序列
\[
0 \to M_s \to M_{n+s} \to M_n \to 0.
\]
它们相容，因为图
\[
\begin{tikzcd}
0 \arrow[r] & M_{s+1} \arrow[r] \arrow[d, "f_{s+1}"] & M_{n+1+s} \arrow[r] \arrow[d, "f_{n+s+1}"] & M_n \arrow[r] \arrow[d, "\mathrm{id}"] & 0 \\
0 \arrow[r] & M_s \arrow[r] & M_{n+s} \arrow[r] & M_n \arrow[r] & 0
\end{tikzcd}
\]
交换.对每个 $n$ 形成上同调序列并对所有 $n$ 取逆向极限，得到正合序列
\[
\cdots \to H^{r}(M) \xrightarrow{\ell^{n}} H^{r}(M) \to H^{r}(M_n) \to H^{r+1}(M) \to \cdots.
\]
这给出正合序列
\[
0 \to H^{r}(M)/\ell^{n}H^{r}(M) \to H^{r}(M_n) \to H^{r+1}(M)[\ell^{n}] \to 0.
\]
因为 $H^{r}(M)$ 是 $\ell$ 幂挠有限群的逆向极限，它没有非零元素被 $\ell$ 的所有幂整除.于是 $\varprojlim H^{r+1}(M)[\ell^{n}] = 0$(转移映射是 $H^{r+1}(M)[\ell^{n}] \xrightarrow{\ell} H^{r+1}(M)[\ell^{n+1}]$)且 $\varprojlim H^{r}(M)/\ell^{n}H^{r}(M) \cong H^{r}(M)$.由此推出 $H^{r}(M)$ 由任何模 $\ell$ 生成它的子集生成.
\end{proof}

\begin{remark}\label{rem:19.3}
紧 $\mathbb{Z}_{\ell}$-模不必有限生成---考虑例如无限个 $\mathbb{F}_{\ell}$ 副本的乘积.

$\mathbb{Z}_{\ell}$-模层 $M = (M_n)$ 称为局部常值的，如果每个 $M_n$ 局部常值.给出连通簇 $X$ 上的局部常值层 $M$ 就是给出有限生成 $\mathbb{Z}_{\ell}$-模 $M$ 连同 $\pi_1(X, \bar{x})$ 在 $M$ 上的连续作用($M$ 上的 $\ell$ 进拓扑).则
\[
H^{1}(X_{\mathrm{et}}, M) = H^{1}_{\mathrm{conts}}(\pi_1(X, \bar{x}), M),
\]
其中 $H^{1}_{\mathrm{conts}}(\pi_1(X, \bar{x}), M)$ 由交叉同态 $f \colon \pi_1(X, \bar{x}) \to M$($M$ 上的 $\ell$ 进拓扑连续的)的等价类组成.

注意，$\mathbb{Z}_{\ell}$-模的局部常值层不必在 $X$ 的任何平展覆盖上变成平凡的，即它不是局部常值的.要使它如此，$\pi_1(X, \bar{x})$ 的作用必须穿过有限商.因此，``$\mathbb{Z}_{\ell}$-模的局部常值层''这个术语是术语的滥用.通常它们被称为``constant-tordu'',``lisse''或``光滑''.
\end{remark}

\subsection*{\texorpdfstring{$\mathbb{Q}_{\ell}$ 模层}{Q\_l 模层}}

$\mathbb{Q}_{\ell}$-向量空间层就是 $\mathbb{Z}_{\ell}$-层 $M = (M_n)$，只是我们定义
\[
H^{r}(X_{\mathrm{et}}, M) = \left( \varprojlim H^{r}(X_{\mathrm{et}}, M_n) \right) \otimes_{\mathbb{Z}_{\ell}} \mathbb{Q}_{\ell}.
\]
例如，
\[
H^{r}(X_{\mathrm{et}}, \mathbb{Q}_{\ell}) = \left( \varprojlim H^{r}(X_{\mathrm{et}}, \mathbb{Z}/\ell^{n}\mathbb{Z}) \right) \otimes_{\mathbb{Z}_{\ell}} \mathbb{Q}_{\ell} = H^{r}(X_{\mathrm{et}}, \mathbb{Z}_{\ell}) \otimes \mathbb{Q}_{\ell}.
\]
在讲义余下部分，$H^{r}(X, \mathbb{Z}_{\ell})$ 表示 $\varprojlim H^{r}(X_{\mathrm{et}}, \mathbb{Z}/\ell^{n}\mathbb{Z})$，$H^{r}(X_{\mathrm{et}}, \mathbb{Q}_{\ell})$ 表示 $H^{r}(X_{\mathrm{et}}, \mathbb{Z}_{\ell}) \otimes \mathbb{Q}_{\ell}$.

\begin{quote}
\textbf{注记} 关于更一般的有限性定理，见 Gabber 的讲义，``Finiteness theorems for \'etale cohomology of excellent schemes''(\href{http://www.math.polytechnique.fr/~laszlo/gdtgabber/abelien.pdf}{此处})与``A finiteness theorem for non abelian $H^{1}$ of excellent schemes''(\href{http://www.math.polytechnique.fr/~laszlo/gdtgabber/non-abelien.pdf}{此处})，以及 Illusie 等人的讲义，``Travaux de Gabber sur l'uniformisation et la cohomologie \'etale des sch\'emas quasi-excellents''(\href{http://www.math.polytechnique.fr/~orgogozo/travaux_de_Gabber/GTG/GTG.pdf}{此处}).
\end{quote}
